% ============================================================
% Notes: Fractional Calculus + Definitions of Fractional Derivatives
% Source: Richard Herrmann, Fractional Calculus: An Introduction for Physicists
% ============================================================

\section{Fractional calculus}

\subsection{Core idea}
Fractional calculus extends differentiation/integration from integer order $n\in\mathbb{N}$
to non-integer order $\alpha\in\mathbb{R}$ (and sometimes complex $\alpha$):
\begin{equation}
\frac{d^n}{dx^n} \quad\longrightarrow\quad \frac{d^\alpha}{dx^\alpha}.
\end{equation}
This provides a systematic framework for \textbf{nonlocal operators} and
\textbf{power-law memory} in physics. :contentReference[oaicite:2]{index=2}

\subsection{Gamma function (replaces factorials)}
Fractional calculus heavily uses the Gamma function:
\begin{equation}
\Gamma(z)=\int_0^\infty t^{z-1}e^{-t}\,dt,
\qquad \Gamma(n+1)=n!.
\end{equation}

% ------------------------------------------------------------

\section{Fractional integrals (Riemann--Liouville)}

\subsection{Riemann--Liouville fractional integral}
For $\alpha>0$, define the fractional integral:
\begin{equation}
\boxed{
(I_{0+}^\alpha f)(t)
=
\frac{1}{\Gamma(\alpha)}\int_0^t (t-\tau)^{\alpha-1}f(\tau)\,d\tau.
}
\end{equation}
This is a Volterra convolution with a power-law kernel, hence it is \textbf{nonlocal in time}.

\subsection{Semigroup property}
For suitable $f$,
\begin{equation}
I^\alpha I^\beta f = I^{\alpha+\beta}f.
\end{equation}

% ------------------------------------------------------------

\section{Fractional derivatives: main definitions}

There is no unique definition of a fractional derivative.
Different definitions agree on broad classes of functions but differ in:
\begin{itemize}
\item how initial conditions are handled,
\item whether constants differentiate to zero,
\item regularity requirements.
\end{itemize}

\subsection{Riemann--Liouville (RL) fractional derivative}
For $n-1<\alpha<n$ with $n\in\mathbb{N}$:
\begin{equation}
\boxed{
(D_{0+}^\alpha f)(t)
=
\frac{d^n}{dt^n}\left(I_{0+}^{\,n-\alpha}f\right)(t)
=
\frac{1}{\Gamma(n-\alpha)}\frac{d^n}{dt^n}
\int_0^t (t-\tau)^{n-\alpha-1}f(\tau)\,d\tau.
}
\end{equation}

\paragraph{Key feature.}
RL derivative is naturally tied to fractional integrals and is mathematically clean,
but it gives (typically)
\begin{equation}
D_{0+}^\alpha (1)\neq 0.
\end{equation}
This can be inconvenient in physics initial-value problems.

% ------------------------------------------------------------

\subsection{Caputo fractional derivative}
For $n-1<\alpha<n$:
\begin{equation}
\boxed{
({}^{C}D_{0+}^\alpha f)(t)
=
I_{0+}^{\,n-\alpha}\left(\frac{d^n f}{dt^n}\right)(t)
=
\frac{1}{\Gamma(n-\alpha)}
\int_0^t (t-\tau)^{n-\alpha-1}f^{(n)}(\tau)\,d\tau.
}
\end{equation}

For $0<\alpha<1$ this simplifies to
\begin{equation}
{}^{C}D_{0+}^\alpha f(t)
=
\frac{1}{\Gamma(1-\alpha)}
\int_0^t (t-\tau)^{-\alpha}f'(\tau)\,d\tau.
\end{equation}

\paragraph{Key feature (physics advantage).}
Caputo derivative satisfies
\begin{equation}
{}^{C}D_{0+}^\alpha (1)=0,
\end{equation}
and initial conditions appear in the usual form $f(0),f'(0),\dots$.
Herrmann emphasizes this as a major motivation for Caputo-type derivatives. :contentReference[oaicite:3]{index=3}

% ------------------------------------------------------------

\subsection{Grünwald--Letnikov (GL) fractional derivative}
A definition closer to the discrete difference quotient:
\begin{equation}
\boxed{
(D^\alpha f)(t)
=
\lim_{h\to 0^+}\frac{1}{h^\alpha}
\sum_{k=0}^{\lfloor t/h\rfloor}(-1)^k\binom{\alpha}{k}\,f(t-kh),
}
\end{equation}
where the generalized binomial coefficient is
\begin{equation}
\binom{\alpha}{k}=\frac{\Gamma(\alpha+1)}{\Gamma(k+1)\Gamma(\alpha-k+1)}.
\end{equation}

\paragraph{Use.}
GL is useful for numerical discretizations and shows explicitly the \textbf{memory}
structure: $f(t)$ depends on all past values $f(t-kh)$.

% ------------------------------------------------------------

\subsection{Fourier/Laplace definitions (spectral characterization)}
Fractional derivatives can be defined via transforms.

\paragraph{Fourier-space (spatial derivative).}
If $\mathcal{F}\{f(x)\}(k)=\tilde f(k)$, define
\begin{equation}
\mathcal{F}\left\{\frac{d^\alpha f}{dx^\alpha}\right\}(k)
=
(ik)^\alpha\,\tilde f(k),
\end{equation}
with a suitable branch choice for $(ik)^\alpha$.

\paragraph{Laplace-space (time derivative).}
For Caputo derivative ($0<\alpha<1$),
\begin{equation}
\boxed{
\mathcal{L}\left\{{}^{C}D_t^\alpha f(t)\right\}(s)
=
s^\alpha F(s) - s^{\alpha-1}f(0).
}
\end{equation}
This is one of the most useful formulas in solving fractional ODEs.

% ------------------------------------------------------------

\section{Herrmann's ``historical'' motivation / consistency conditions}

Herrmann motivates fractional derivatives by demanding that standard derivative rules
extend smoothly to non-integer order for simple eigenfunctions. :contentReference[oaicite:4]{index=4}

\subsection{Exponentials as eigenfunctions}
\begin{equation}
\boxed{
\frac{d^\alpha}{dx^\alpha}e^{kx} = k^\alpha e^{kx},
\qquad k\ge 0.
}
\end{equation}

\subsection{Trigonometric functions}
\begin{equation}
\boxed{
\frac{d^\alpha}{dx^\alpha}\sin(kx)
=
k^\alpha\sin\!\left(kx+\frac{\pi\alpha}{2}\right),
\qquad k\ge 0.
}
\end{equation}

\subsection{Power functions}
For $x\ge 0$,
\begin{equation}
\boxed{
\frac{d^\alpha}{dx^\alpha}x^m
=
\frac{\Gamma(m+1)}{\Gamma(m+1-\alpha)}x^{m-\alpha}.
}
\end{equation}

These identities are consistent with integer derivatives when $\alpha=n\in\mathbb{N}$.

% ------------------------------------------------------------

\section{RL vs Caputo (quick comparison)}

For $0<\alpha<1$:
\begin{itemize}
\item RL:
\[
D^\alpha f(t)=\frac{d}{dt}\left[\frac{1}{\Gamma(1-\alpha)}\int_0^t (t-\tau)^{-\alpha}f(\tau)\,d\tau\right].
\]
\item Caputo:
\[
{}^{C}D^\alpha f(t)=\frac{1}{\Gamma(1-\alpha)}\int_0^t (t-\tau)^{-\alpha}f'(\tau)\,d\tau.
\]
\end{itemize}

Relation (for sufficiently smooth $f$):
\begin{equation}
{}^{C}D^\alpha f(t)
=
D^\alpha\big(f(t)-f(0)\big).
\end{equation}

% ------------------------------------------------------------

\section{High-yield summary}

\begin{itemize}
\item Fractional calculus extends derivatives/integrals to order $\alpha\notin\mathbb{N}$.
\item RL fractional integral:
\[
(I^\alpha f)(t)=\frac{1}{\Gamma(\alpha)}\int_0^t (t-\tau)^{\alpha-1}f(\tau)\,d\tau.
\]
\item RL derivative:
\[
D^\alpha f=\frac{d^n}{dt^n}I^{n-\alpha}f,\quad n-1<\alpha<n.
\]
\item Caputo derivative:
\[
{}^{C}D^\alpha f=I^{n-\alpha}f^{(n)},\quad n-1<\alpha<n,
\]
and ${}^{C}D^\alpha(\text{const})=0$.
\item GL derivative gives a discrete memory sum and is useful numerically.
\item Transform rule:
\[
\mathcal{L}\{{}^{C}D_t^\alpha f\}=s^\alpha F(s)-s^{\alpha-1}f(0).
\]
\item Herrmann motivation: preserve eigenfunction rules
(e.g.\ $d^\alpha e^{kx}/dx^\alpha=k^\alpha e^{kx}$). :contentReference[oaicite:5]{index=5}
\end{itemize}
